\section{Conclusion}
\label{sec:conclusion}

Coding agents such as Claude Code turn the language model from an advisor into an operator, and
the human from operator into approver or observer \citep{feng2025levels}. This report examined one
risk of that transition: that the sub-tasks an agent meets or receives during a session---broken
dependencies, failing tests, instructions in files, side requests---accumulate and pull it away from
the goal it was given. We reviewed the work on autonomy, coding agents, oversight and goal drift. We
defined goal drift via sub-task injection as a counterfactual degradation, with sub-tasks classified
by source and legitimacy and a matched-length control. We then tested the definition on
DriftBench-mini, with Claude Code driven by two open-weight models in 1{,}332 controlled episodes and a 72-episode compaction pilot.

The findings revise the starting intuition:
\begin{itemize}
  \item \textbf{Kind matters more than count.} Neutral interruptions and mild distractors left goal
        completion unchanged. Broken dependencies on the path to the goal and, for the weaker model,
        a stack of four side requests reduced it sharply.
  \item \textbf{The failure is incompletion.} The agents rarely strayed outside the files they were
        supposed to change. They lost the goal by not finishing it, in a way that a diff-based review
        would not reveal.
  \item \textbf{Lightweight harness mitigations are not enough.} Re-injecting the goal, a stop-hook
        checklist, an instruction to delegate and an enforced task list did not significantly reduce
        the loss. The two that relied on the model's compliance were adopted only partly.
  \item \textbf{Off-goal work is mostly legitimate.} A panel of three LLM judges, reliable on the
        binary off-goal label ($\kappa=0.67$) but not validated against humans, found that the work
        leaving the goal was almost all what the user asked for or what the goal needed. Work on
        distractors stayed below 7\% of segments in every cell. Final summaries, however, claimed
        completion with a hidden test failing in 42--50\% of baseline episodes.
  \item \textbf{Compaction may tilt effort towards the requests.} In a forced-compaction pilot,
        \texttt{qwen3-coder} spent more of the segments after compaction on the user's side requests
        (off-goal share $+0.13$), with a small and uncertain loss on the goal. The pilot is small and
        its placebo not clean.
  \item \textbf{Watching the activity catches few failures.} In a simulated supervisor with perfect
        labels, reviewing every segment flagged about half of the failing episodes before the end at
        six reviews per episode and with frequent false alarms, and a monitor for illegitimate work
        flagged fewer than one in five. Failing episodes look like goal work; only the requirements
        show the failure.
\end{itemize}
For the human side of the interaction, these results argue for oversight that checks requirements
instead of code, checkpoints triggered at recognisable risky moments instead of fixed intervals,
mid-session requests handled as queued tasks, and mitigations that the harness enforces without
relying on the model's compliance.

Four lines of work follow from what this study could not settle:
\begin{itemize}
  \item \emph{inherited drift}, with a weak model that uses the harness competently, so that its
        prefixes drift instead of breaking;
  \item \emph{monitor degradation}: whether a drift monitor loses accuracy on longer transcripts
        \citep{martin2026classifier}, which decides whether drift-triggered checkpoints scale;
  \item \emph{human validation} of the judge panel on a sample, since its agreement shows
        reliability but not validity;
  \item \emph{goal recall}, on which the judges do not yet agree ($\kappa=0.17$), and which is needed
        to test whether compaction changes what the agent believes the goal to be.
\end{itemize}
Repeating the core experiments with frontier models, longer and less structured sessions, and
real-world rather than planted sub-tasks would show how far the thresholds observed here move as
agents become more capable. We expect them to move; our data cannot say whether they disappear.
